\documentclass[11pt,a4paper]{article}

% -----------------------------------------------------------------------------
% CS6868 Research Mini-Project Report
% Project 23: Synchronous Dual Queue
% -----------------------------------------------------------------------------

\usepackage[utf8]{inputenc}
\usepackage[T1]{fontenc}
\usepackage[margin=1in]{geometry}
\usepackage{microtype}
\usepackage{lmodern}
\usepackage{amsmath,amssymb}
\usepackage{graphicx}
\usepackage{booktabs}
\usepackage{array}
\usepackage{enumitem}
\usepackage{xcolor}
\usepackage{listings}
\usepackage{hyperref}
\usepackage[capitalise,noabbrev]{cleveref}

\hypersetup{
  colorlinks=true,
  linkcolor=black,
  citecolor=blue!60!black,
  urlcolor=blue!60!black,
}

% --- Listings style (for code snippets) ---------------------------------------
\definecolor{codegray}{gray}{0.95}
\lstset{
  basicstyle=\ttfamily\small,
  backgroundcolor=\color{codegray},
  frame=single,
  framesep=4pt,
  breaklines=true,
  columns=fullflexible,
  keepspaces=true,
  showstringspaces=false,
}

\lstdefinelanguage{caml}{
  keywords={let,rec,in,match,with,fun,function,if,then,else,begin,end,type,of,module,struct,sig,val,and,Some,None,true,false,ref},
  sensitive=true,
  comment=[s]{(*}{*)},
  morestring=[b]"
}

% --- Title block --------------------------------------------------------------
\title{%
  \textbf{CS6868 Research Mini-Project} \\[0.3em]
  \Large Project 23: A Lock-Free Synchronous Dual Queue in OCaml 5%
}
\author{%
  Rohan Kadam \\ \texttt{cs22b083@smail.iitm.ac.in}
  \and
  Bhargav \\ \texttt{cs22b087@smail.iitm.ac.in}
  \and
  Venkat \\ \texttt{cs22b027@smail.iitm.ac.in}
}
\date{\today}

\begin{document}
\maketitle

\begin{abstract}
We implemented a lock-free synchronous dual queue based on the design of
Scherer, Lea and Scott~\cite{scherer2009synchronous}, together with two
baselines: a blocking mutex/condition-variable synchronous queue, and a
single-slot lock-free exchanger from Lecture~08. All three were written in
OCaml~5 using the new \texttt{Atomic} module and \texttt{Domain} parallelism.
We verified correctness using manual concurrent stress tests and a
QCheck-Lin linearizability test on the blocking baseline, and we ran the
lock-free implementation under ThreadSanitizer (TSAN) on a TSAN-enabled
OCaml switch. Throughput was measured in matched pairs per second across
2--8 domains under balanced, enqueue-heavy and dequeue-heavy workloads.
The lock-free dual queue beats the blocking baseline by roughly
$4$--$10\times$ across every tested configuration, including the
imbalanced workloads, so the reservation-node design is a clear win
over mutex/condition variables at our scales. The exchanger, which
only handles paired threads in the balanced configuration, adds a
further $1.2$--$2.9\times$ on top of the dual queue and, unlike the
queue, its throughput scales near-linearly with the domain count.
This suggests that elimination on paired arrivals is a worthwhile
optimisation to layer on top of the dual queue, which we discuss as
future work.
\end{abstract}

% =============================================================================
\section{Goals}
\label{sec:goals}
% =============================================================================

We picked Project~23 (Synchronous Dual Queue) because the algorithm
neatly combines two ideas the course covered separately: the
Michael--Scott lock-free queue and the use of CAS
to coordinate handoff between threads (Lecture~08, exchanger and
elimination array). A standard concurrent queue is asynchronous: an
enqueuer always succeeds, even if no dequeuer is around. A
\emph{synchronous} queue, in contrast, requires a rendezvous: an
\texttt{enqueue} must wait until some \texttt{dequeue} actually receives
the value, and vice versa. Java's \texttt{SynchronousQueue} and Go
channels of capacity zero are real systems that work this way. The
\emph{dual queue} of Scherer, Lea and Scott~\cite{scherer2009synchronous}
unifies both modes by allowing reservation nodes to live in the same
Michael--Scott structure, with producers fulfilling consumer reservations
and vice versa.

The connection to course material is direct: the dual queue is the
Michael--Scott queue from Lecture~08 plus condition synchronization done
through CAS rather than condition variables, which is exactly what
Scherer and Scott call ``nonblocking concurrent data structures
with condition synchronization.''

\paragraph{Research question.}
\emph{How does the lock-free dual queue's rendezvous throughput compare
against a mutex/condition-based synchronous queue, and does the overhead
of reservation nodes pay off at higher thread counts?}

% =============================================================================
\section{Background}
\label{sec:background}
% =============================================================================

\paragraph{Synchronous queues.}
A synchronous queue has two operations, \texttt{put} and \texttt{take},
both of which block until they meet a partner of the opposite kind.
Conceptually it is a buffer of capacity zero: every value is handed off
directly from a producer to a consumer with no storage in between. The
classical implementation in textbooks~\cite[\S10.6.1]{AoMPP} uses a
mutex and two condition variables, one for waiting producers and one
for waiting consumers, and every \texttt{put}/\texttt{take} pair causes
two context switches.

\paragraph{Dual queue intuition.}
The Scherer--Lea--Scott dual queue avoids the lock by storing the
``waiting'' state inside the queue itself. The queue holds nodes of two
kinds:
\begin{itemize}[leftmargin=1.4em,itemsep=0.2em]
  \item \texttt{Data}: a producer is waiting; the node carries
    \texttt{Some v}.
  \item \texttt{Request}: a consumer is waiting; the node's value slot
    is initially \texttt{None}.
\end{itemize}
The key invariant is that all live nodes in the queue have the
\emph{same} kind --- the queue is either empty, all-Data, or all-Request,
never mixed. An arriving thread inspects the head:
\begin{itemize}[leftmargin=1.4em,itemsep=0.2em]
  \item If the queue is empty or only contains waiters of the same kind
    as the arriving thread, the thread appends a reservation (using a
    Michael--Scott-style two-step tail CAS) and spins until its own node
    is fulfilled.
  \item Otherwise, the head holds an opposite-kind reservation; the
    arriving thread \emph{fulfils} that reservation by CAS-ing the
    payload slot from its waiting state to its fulfilled state, then
    helps advance the head.
\end{itemize}
The CAS that flips the payload slot is the rendezvous: it is the
linearization point for both the fulfilling thread and the waiting
thread, simultaneously. After the CAS, the waiter sees its node is
fulfilled and returns; the fulfiller helps advance \texttt{head} and
also returns. \Cref{lst:put} sketches the producer side.

\begin{lstlisting}[language=caml,caption={Producer (\texttt{put}) sketch.},label={lst:put}]
let rec put q value =
  help_advance_head q;
  match Atomic.get (Atomic.get q.head).next with
  | Some ({ kind = Request; item; _ } as r) ->
      (* Consumer is waiting: fulfil the head request. *)
      if Atomic.compare_and_set item None (Some value) then
        ignore (Atomic.compare_and_set q.head ... r)
      else put q value
  | _ ->
      (* Empty or only producers waiting: append reservation and spin. *)
      let node = {kind=Data; item=Atomic.make (Some value); next=Atomic.make None} in
      match try_append_node q node with
      | Appended -> wait_until_fulfilled node
      | Retry    -> put q value
\end{lstlisting}

\paragraph{Linearizability.}
A history is linearizable if every
operation appears to take effect at a single instant between its
invocation and response. For our dual queue, that instant is the
successful payload CAS at the rendezvous. Both \texttt{put} and
\texttt{take} share the same linearization point: this is the elegant
part of the algorithm and the part that took us the longest to
internalise.

% =============================================================================
\section{Tasks Undertaken}
\label{sec:tasks}
% =============================================================================

\subsection{Implementation}

We implemented three handoff primitives, all in OCaml~5 using
\texttt{Atomic.t} and \texttt{Domain.spawn}. Source paths below are
relative to the repository root.

\begin{itemize}[leftmargin=1.4em,itemsep=0.2em]
  \item \textbf{Lock-free synchronous dual queue}
    (\texttt{code/lockfree\_sync\_dual\_queue.ml}, 166 LoC). Nodes carry
    an \texttt{'a option Atomic.t} payload so the same atomic flips
    represent both the value and the fulfilment state. We use
    \texttt{Atomic.make\_contended} for \texttt{head}, \texttt{tail} and
    every node's \texttt{item} slot to push them onto separate cache
    lines and reduce false sharing. \texttt{help\_advance\_head} is
    called eagerly at the top of \texttt{put}/\texttt{take} so that
    fulfilled nodes do not pile up at the front.
  \item \textbf{Single-slot exchanger}
    (\texttt{code/lockfree\_exchanger.ml}, 61 LoC). One
    \texttt{state Atomic.t} cycles through \texttt{Empty}, \texttt{Waiting v},
    \texttt{Busy v}. The first arrival CASes \texttt{Empty $\to$ Waiting v}
    and spins; the second arrival CASes \texttt{Waiting v $\to$ Busy v'}
    and returns. The successful second CAS is the linearization point
    for both threads.
  \item \textbf{Mutex + condition baseline}
    (\texttt{code/lock\_sync\_dual\_queue.ml}, 107 LoC). A single mutex
    protects two \texttt{Stdlib.Queue} structures of waiting producers
    and waiting consumers, each carrying its own
    \texttt{Condition.t}. We also expose non-blocking \texttt{try\_put}
    and \texttt{try\_take} variants that we use for QCheck-Lin testing
    (the blocking versions cause QCheck workers to deadlock).
\end{itemize}

\paragraph{Subtle invariants.}
The dual queue has an invariant we got wrong twice during development:
the queue may only contain reservations \emph{of one kind at a time}.
The first version of \texttt{put} appended a Data node whenever it saw
a non-empty queue, which broke the invariant on a tail-CAS race and
caused two consumers to mis-pair with the same producer in stress
tests. Our fix is the conjunctive guard
\begin{lstlisting}[language=caml]
if last == first || last.kind = node.kind then ...
\end{lstlisting}
which mirrors how the original paper splits the head-fulfil and
tail-append cases. We also re-read \texttt{q.tail} after the snapshot
to detect concurrent appends, in the same way Michael--Scott
re-validates \texttt{tail}.

\paragraph{Two-step tail.}
We deliberately kept the two-step tail update (set \texttt{last.next}
first, then CAS \texttt{q.tail}) because it lets a stalled appender
hand off the rest of its work to the next thread, which is essential
for lock-freedom in OCaml's runtime where any \texttt{Domain} can be
descheduled.

\subsection{Testing and verification}

\paragraph{Manual stress tests
(\texttt{code/test\_manual.ml}).}
We wrote five concurrent tests, each exercising a different
configuration: 1p/1c single pair, 1p/1c repeated pairs (10 and 100),
4p/4c with 200 pairs and a per-value duplicate detector, plus exchanger
tests for 2, 4 and 8 domains. The 4p/4c test was the most useful: each
producer publishes values $\{i, i+P, i+2P, \dots\}$ where $P$ is the
producer count, the consumers record receipts in a shared array
guarded by a mutex, and the test fails if any value is missed,
duplicated or out of range. This caught the same-kind invariant bug
mentioned above.

\paragraph{QCheck-Lin
(\texttt{code/qcheck\_lin.ml}).}
We use QCheck-Lin's \texttt{Lin\_domain} backend with the
\texttt{multicoretests} library. Because the
blocking variants of \texttt{put}/\texttt{take} would simply deadlock
the QCheck-Lin workers (there is no consumer when QCheck-Lin tries to
linearise a single producer's history), we expose
\texttt{try\_put}/\texttt{try\_take} on the lock baseline and
exercise those instead. QCheck-Lin runs 500 random concurrent
histories per invocation and checks that each one is sequentially
consistent against an explicit ``either rendezvous now or fail''
specification. All 500 traces passed.

\paragraph{TSAN.}
We built the project on a TSAN-enabled OCaml switch
(an \texttt{options+tsan} variant of OCaml~5) and re-ran
\texttt{test\_manual} and a 1\,000-pair benchmark
(\texttt{results/benchmark-tsan.csv}). TSAN reported no data races on
the lock-free queue. As a sanity check we also wrote a deliberately
broken non-atomic version
(\texttt{code/non\_atomic\_lockfree\_sync\_dual\_queue\_tsan.ml}) that
uses plain \texttt{ref} cells for the same fields, and confirmed that
TSAN does flag races there. This convinced us that the TSAN setup
was actually instrumenting the code rather than silently passing.

\paragraph{Producer--consumer pipeline
(\texttt{code/producer\_consumer\_pipeline/pipeline\_demo.ml}).}
A small demo with 4 producers, 8 consumers and 10 jobs per producer.
The main thread spawns all consumers, spawns all producers, joins the
producers, and then sends one \texttt{Stop} message per consumer through
the same synchronous queue. Because every handoff is a rendezvous,
the demo verifies that producers and consumers actually wait for each
other rather than buffering through some implicit queue.

% =============================================================================
\section{Evaluation}
\label{sec:evaluation}
% =============================================================================

\paragraph{Setup.}
All measurements were taken on a Linux laptop (kernel 6.17.0-22-generic,
x86\_64, 12th Gen Intel Core i5-1235U, 10 cores / 12 logical threads,
OCaml~5.4.0). The benchmark
(\texttt{code/benchmark.ml}) spawns the requested number of producer
and consumer domains, has them all spin on a shared
\texttt{Atomic.bool} ``go'' flag, and times the wall-clock interval
from setting the flag to all domains joining. Each producer
\texttt{put}s a fixed share of the total pair count, each consumer
\texttt{take}s its share, and we report
$\text{pairs}/\text{seconds}$ as throughput. We use 100\,000 pairs per
configuration, which was enough to dominate the start-up jitter on our
machine. Because of the tight time budget for the project we did not
sweep over multiple repetitions and did not pin domains; the trends are
robust but the absolute numbers carry $\pm$5--10\% run-to-run noise,
which is something we would tighten given more time.

\paragraph{Workloads.}
\emph{Balanced} splits domains evenly between producers and
consumers; \emph{enqueue-heavy} uses $\lceil 3D/4 \rceil$ producers
and the rest consumers; \emph{dequeue-heavy} is the mirror image.
Every configuration moves the same total of 100\,000 matched pairs, so
imbalance changes the ratio of waiters to fulfillers but not the total
amount of synchronisation work.

\begin{table}[h]
  \centering
  \small
  \caption{Throughput (pairs/s, $\times10^3$) at 100\,000 pairs.
  Source: \texttt{code/results/benchmark.csv}.}
  \label{tab:throughput}
  \begin{tabular}{@{}llrrrr@{}}
    \toprule
    Workload & Implementation & 2D & 4D & 6D & 8D \\
    \midrule
    Balanced       & lockfree  & 1047.1 &  978.8 & 1242.8 & 1203.1 \\
                   & blocking  &  243.6 &  249.0 &  210.2 &  207.3 \\
    \midrule
    Enqueue-heavy  & lockfree  & 1556.9 & 1307.0 & 1313.4 & 1085.9 \\
                   & blocking  &  191.0 &  241.9 &  218.8 &  149.0 \\
    \midrule
    Dequeue-heavy  & lockfree  & 1461.2 & 1415.2 &  972.5 & 1069.4 \\
                   & blocking  &  146.4 &  235.4 &  157.0 &  158.9 \\
    \bottomrule
  \end{tabular}
\end{table}

\paragraph{Balanced workload.}
On the balanced workload (\cref{tab:throughput}) the lock-free dual
queue is $4.3\times$ faster than the blocking baseline at 2 domains,
$3.9\times$ at 4, $5.9\times$ at 6, and $5.8\times$ at 8. The gap is
large even at 2 domains because the mutex path pays for a real
\texttt{pthread\_cond\_wait}/\texttt{signal} pair on every handoff,
whereas the lock-free rendezvous is just two CASes and a short
\texttt{cpu\_relax} spin. Throughput on the lock-free queue is
roughly flat in the 1.0--1.2\,M pairs/s range across domain counts,
which is consistent with the tail (and the head) being the serial
bottleneck --- adding more domains does not add more handoff
parallelism when every pair still funnels through the same
\texttt{head}/\texttt{tail} atomics.

\paragraph{Imbalanced workloads.}
The enqueue-heavy and dequeue-heavy columns show the same pattern:
the lock-free queue stays in the 1.0--1.6\,M pairs/s band, while the
blocking baseline drops to 0.15--0.25\,M. The lock-free advantage
widens to $5$--$10\times$ under imbalance. The reason is that the
lock-free fast path is the same regardless of who is waiting: a
fulfiller arrives, CASes the head reservation's payload slot, and
leaves. The blocking implementation, by contrast, has to take the
mutex, push the waiter onto the appropriate queue, release the
mutex, \texttt{wait} on a condition variable, and then wake up and
re-acquire when signalled --- and the more often one side is
waiting (which is exactly what ``imbalanced'' means), the more
often that full wait/signal cycle runs. The 2-domain
dequeue-heavy point for the blocking baseline (146\,k pairs/s) is
the worst of the blocking numbers, again consistent with this: a
single consumer repeatedly parks on the empty queue, and every
producer has to wake it.

\paragraph{The exchanger.}
The exchanger is faster than the dual queue on the balanced
workload but not by as much as one might expect ---
$1.18\times$ at 2 domains, $2.35\times$ at 4, $2.90\times$ at 6,
and $2.69\times$ at 8. Unlike the dual queue, the exchanger's
throughput \emph{does} scale with the domain count, because each
pair has its own isolated \texttt{state Atomic.t} (in our harness
we use a single slot, so the scaling we see is really
parallelism in the two producer/consumer halves rather than
across slots); with a real elimination array of $D$ slots we
would expect the gap to the queue to widen further. Either way,
the fast path allocates nothing and consists of two CASes, so it
is a strict lower bound on what any handoff primitive can cost.
We see this as motivation for adding an elimination layer on top
of the dual queue, similar to the elimination-backoff stack from
Lecture~08: arrive at an exchanger slot first, only fall back to
the slow Michael--Scott path if there is no partner within a
budget.

\paragraph{Did the reservation overhead pay off?}
Yes, unambiguously. Across all twelve queue configurations
($\{$balanced, enqueue-heavy, dequeue-heavy$\} \times \{2,4,6,8\}$
domains) the lock-free dual queue beats the mutex/condition-variable
baseline by no less than $3.9\times$ and up to $10.0\times$. The
reservation-node cost (one extra allocation per waiting thread, plus
a few atomic loads to find the queue kind) is much smaller than the
cost of a kernel-level park/unpark, even on a single-machine
workload with no real I/O.

% =============================================================================
\section{Reflection on the Use of LLMs}
\label{sec:llm}
% =============================================================================

We used Cursor and ChatGPT while working on this project.

% =============================================================================
\section{Conclusions}
\label{sec:conclusions}
% =============================================================================

Our research question was: \emph{does the overhead of reservation
nodes pay off at higher thread counts?} The answer we got is
``yes, by a wide margin.'' Across balanced, enqueue-heavy and
dequeue-heavy workloads at 2--8 domains, the lock-free dual queue
gives $3.9$--$10.0\times$ the throughput of the mutex/condition
baseline. The win is driven by the two fast paths that the
reservation design enables: an arriving thread that meets an
opposite-kind waiter finishes the rendezvous in two CASes, and an
arriving thread that has to wait spins on its own node rather than
parking in the kernel. Both paths avoid the \texttt{pthread\_cond}
wake-up that dominates the blocking baseline's cost. The only
implementation whose throughput \emph{grows} with the domain count
in our data is the exchanger, which hints that the queue's shared
\texttt{head}/\texttt{tail} atomics are now the bottleneck.

\paragraph{Limitations.}
Our benchmark is a single machine, single OCaml version, no domain
pinning, and one repetition per point; the 5--10\% run-to-run noise
we observed is enough that some of the smaller differences in
\cref{tab:throughput} (for example lockfree at 2D vs.\ 4D on
balanced) should be read as ``no measurable difference.'' We also
did not implement the full elimination-array variant of the dual
queue from Scherer--Lea--Scott, which is what would close the
remaining $\sim$2--3$\times$ gap to the single-slot exchanger.

\paragraph{Future work.}
The two natural next steps are (i) layering the single-slot
exchanger (or an elimination array of them) in front of the slow
path to capture paired arrivals without queueing, and (ii)
replacing our unbounded spin with bounded spin + \texttt{Domain}
parking, which would let the dual queue degrade gracefully when the
number of waiters blows past the number of cores. Both are
well-known fixes from the literature and from Java's
\texttt{SynchronousQueue}~\cite{javasynchronousqueue}.

% =============================================================================
\section*{Contributions}
\label{sec:contributions}
% =============================================================================

\begin{center}
\begin{tabular}{@{}lr@{}}
\toprule
\textbf{Member} & \textbf{Contribution (\%)} \\
\midrule
Rohan Kadam   & 34\% \\
Bhargav       & 33\% \\
Venkat        & 33\% \\
\midrule
\textbf{Total} & \textbf{100\%} \\
\bottomrule
\end{tabular}
\end{center}

% =============================================================================
\bibliographystyle{plainurl}
\bibliography{references}
% =============================================================================

\end{document}
